\documentclass[11pt,letterpaper]{article}
\usepackage[technical-reference]{cornell-notes}
\usepackage{needspace}

\title{Clause 8: Statements - Cornell Notes}
\author{}
\date{}
\setDocTitle{Clause 8: Statements}
\setDocSubtitle{C++ 2024 Cornell Notes}
\setDocOwner{C++ 2024 Cornell Notes}
\setDocAuthor{}
\setDocDate{}
\setCornellCollection{C++ 2024 Cornell Notes}
\setCornellUnitType{Clause}
\setCornellUnitNumber{8}
\setCornellUnitTitle{Statements}
\setCornellCompanionLabel{C++ Technical Reference}
\hypersetup{pdftitle={Clause 8: Statements - Cornell Notes},pdfsubject={C++ technical study notes},pdfauthor={}}

\begin{document}
\maketitle
\begin{CornellReferenceOrientationBox}
\textbf{Purpose.} Defines labels, blocks, selection, iteration, jumps, declaration statements, coroutine returns, and grammar ambiguity resolution.
\par\textbf{Role.} Organizes expressions and declarations into structured control flow.
\par\textbf{Principal dependencies.} Uses expressions and declarations; interacts with scope, lifetime, exceptions, and coroutine rules.
\par\textbf{Primary audience.} programmers and control-flow analysis tools.
\par\textbf{Major subject groups.} Preamble; Label; Expression statement; Compound statement or block; Selection statements; Iteration statements; Jump statements; Declaration statement; Ambiguity resolution.
\end{CornellReferenceOrientationBox}
\section{Learning Objectives}
\begin{itemize}[label=\textcolor{CornellReferenceTeal}{$\blacktriangleright$}]
\item Define the governing ideas of Preamble and state the consequences for a program or implementation.
\item Identify the governing ideas of Label and state the consequences for a program or implementation.
\item Distinguish the governing ideas of Expression statement and state the consequences for a program or implementation.
\item Explain the governing ideas of Compound statement or block and state the consequences for a program or implementation.
\item Trace the governing ideas of Selection statements and state the consequences for a program or implementation.
\item Apply the governing ideas of Iteration statements and state the consequences for a program or implementation.
\item Analyze the governing ideas of Jump statements and state the consequences for a program or implementation.
\item Evaluate the governing ideas of Declaration statement and state the consequences for a program or implementation.
\item Diagnose the governing ideas of Ambiguity resolution and state the consequences for a program or implementation.
\item Compare the governing ideas of General and state the consequences for a program or implementation.
\end{itemize}
\section{Normative Structure Map}
\small\begin{longtable}{@{}>{\RaggedRight\arraybackslash}p{0.13\textwidth}>{\RaggedRight\arraybackslash}p{0.27\textwidth}>{\RaggedRight\arraybackslash}p{0.19\textwidth}>{\RaggedRight\arraybackslash}p{0.31\textwidth}@{}}
\toprule\textbf{Subclause} & \textbf{Subject} & \textbf{Rule category} & \textbf{Key dependencies / entities}\\\midrule\endfirsthead
\toprule\textbf{Subclause} & \textbf{Subject} & \textbf{Rule category} & \textbf{Key dependencies / entities}\\\midrule\endhead
\bottomrule\endfoot
8.1 & Preamble & core-language semantics & Uses expressions and declarations; interacts with scope, lifetime, exceptions, and coroutine rules.\\
8.2 & Label & core-language semantics & Uses expressions and declarations; interacts with scope, lifetime, exceptions, and coroutine rules.\\
8.3 & Expression statement & core-language semantics & Uses expressions and declarations; interacts with scope, lifetime, exceptions, and coroutine rules.\\
8.4 & Compound statement or block & concurrency contract & Uses expressions and declarations; interacts with scope, lifetime, exceptions, and coroutine rules.\\
8.5 & Selection statements & core-language semantics & General, The if statement, The switch statement\\
8.6 & Iteration statements & core-language semantics & General, The while statement, The do statement, The for statement, and related subclauses\\
8.7 & Jump statements & core-language semantics & General, The break statement, The continue statement, The return statement, and related subclauses\\
8.8 & Declaration statement & core-language semantics & Uses expressions and declarations; interacts with scope, lifetime, exceptions, and coroutine rules.\\
8.9 & Ambiguity resolution & core-language semantics & Uses expressions and declarations; interacts with scope, lifetime, exceptions, and coroutine rules.\\
\end{longtable}\normalsize
\begin{CornellReferenceNormativeBox}A heading maps the clause; it does not replace the detailed conditions. For each operation, read requirements, effects, result, exceptions, complexity, remarks, and notes with their stated normative force.\end{CornellReferenceNormativeBox}
\subsection*{Rule Classification Guide}
\small\begin{longtable}{@{}>{\RaggedRight\arraybackslash}p{0.25\textwidth}>{\RaggedRight\arraybackslash}p{0.67\textwidth}@{}}\toprule\textbf{Label or category} & \textbf{How to read it}\\\midrule\endfirsthead\toprule\textbf{Label or category} & \textbf{How to read it (continued)}\\\midrule\endhead\bottomrule\endfoot
Well-formed / ill-formed & Formation is governed by syntax and semantic rules; an ill-formed program does not become well-formed because one implementation accepts it.\\
Diagnosable rule & A violation requires at least one diagnostic unless the wording places the case outside the diagnosable set.\\
No diagnostic required & The program can be ill-formed while the implementation has no diagnostic obligation.\\
Undefined behavior & No execution requirements are imposed; translation success does not establish safety or portability.\\
Unspecified behavior & One of the permitted outcomes may be chosen without documentation.\\
Implementation-defined behavior & The implementation chooses and documents the behavior.\\
Conditionally supported & The implementation may omit support but documents unsupported constructs.\\
\end{longtable}\normalsize
\section{Cornell Cue-and-Notes}
\Needspace{8\baselineskip}
\subsection*{8.1 Preamble}
\begin{longtable}{@{}>{\RaggedRight\arraybackslash\columncolor{CornellReferenceCueGray}}p{0.28\textwidth}>{\RaggedRight\arraybackslash}p{0.64\textwidth}@{}}
\rowcolor{CornellReferenceNavy}\color{white}\textbf{Cue / recall prompt} & \color{white}\textbf{Notes}\\\endfirsthead
\rowcolor{CornellReferenceNavy}\color{white}\textbf{Cue / recall prompt} & \color{white}\textbf{Notes (continued)}\\\endhead
\arrayrulecolor{CornellReferenceLineGray}
\textbf{What rule applies to Preamble?}\\[-1mm]{\scriptsize\CornellClauseRef{8.1}} & Frames the terminology, applicability, and relationships needed to interpret Preamble; detailed rules remain in the following subclauses.\\\hline
\end{longtable}
\begin{CornellReferenceSummaryBox}Preamble supplies the core-language semantics portion of this clause. The key review move is to connect its local wording to the clause-wide model, then classify each condition by normative force before relying on it.\end{CornellReferenceSummaryBox}
\Needspace{8\baselineskip}
\subsection*{8.2 Label}
\begin{longtable}{@{}>{\RaggedRight\arraybackslash\columncolor{CornellReferenceCueGray}}p{0.28\textwidth}>{\RaggedRight\arraybackslash}p{0.64\textwidth}@{}}
\rowcolor{CornellReferenceNavy}\color{white}\textbf{Cue / recall prompt} & \color{white}\textbf{Notes}\\\endfirsthead
\rowcolor{CornellReferenceNavy}\color{white}\textbf{Cue / recall prompt} & \color{white}\textbf{Notes (continued)}\\\endhead
\arrayrulecolor{CornellReferenceLineGray}
\textbf{What rule applies to Label?}\\[-1mm]{\scriptsize\CornellClauseRef{8.2}} & Defines the core-language rules for Label, including formation, semantic effect, interactions with type and lookup, and any ill-formed or implementation-dependent cases.\\\hline
\end{longtable}
\begin{CornellReferenceSummaryBox}Label supplies the core-language semantics portion of this clause. The key review move is to connect its local wording to the clause-wide model, then classify each condition by normative force before relying on it.\end{CornellReferenceSummaryBox}
\Needspace{8\baselineskip}
\subsection*{8.3 Expression statement}
\begin{longtable}{@{}>{\RaggedRight\arraybackslash\columncolor{CornellReferenceCueGray}}p{0.28\textwidth}>{\RaggedRight\arraybackslash}p{0.64\textwidth}@{}}
\rowcolor{CornellReferenceNavy}\color{white}\textbf{Cue / recall prompt} & \color{white}\textbf{Notes}\\\endfirsthead
\rowcolor{CornellReferenceNavy}\color{white}\textbf{Cue / recall prompt} & \color{white}\textbf{Notes (continued)}\\\endhead
\arrayrulecolor{CornellReferenceLineGray}
\textbf{What rule applies to Expression statement?}\\[-1mm]{\scriptsize\CornellClauseRef{8.3}} & Defines the core-language rules for Expression statement, including formation, semantic effect, interactions with type and lookup, and any ill-formed or implementation-dependent cases.\\\hline
\end{longtable}
\begin{CornellReferenceSummaryBox}Expression statement supplies the core-language semantics portion of this clause. The key review move is to connect its local wording to the clause-wide model, then classify each condition by normative force before relying on it.\end{CornellReferenceSummaryBox}
\Needspace{8\baselineskip}
\subsection*{8.4 Compound statement or block}
\begin{longtable}{@{}>{\RaggedRight\arraybackslash\columncolor{CornellReferenceCueGray}}p{0.28\textwidth}>{\RaggedRight\arraybackslash}p{0.64\textwidth}@{}}
\rowcolor{CornellReferenceNavy}\color{white}\textbf{Cue / recall prompt} & \color{white}\textbf{Notes}\\\endfirsthead
\rowcolor{CornellReferenceNavy}\color{white}\textbf{Cue / recall prompt} & \color{white}\textbf{Notes (continued)}\\\endhead
\arrayrulecolor{CornellReferenceLineGray}
\textbf{What rule applies to Compound statement or block?}\\[-1mm]{\scriptsize\CornellClauseRef{8.4}} & Defines the blocking and synchronization contract for Compound statement or block; ownership, wake conditions, timeout behavior, and happens-before edges must be analyzed together.\\\hline
\end{longtable}
\begin{CornellReferenceSummaryBox}Compound statement or block supplies the concurrency contract portion of this clause. The key review move is to connect its local wording to the clause-wide model, then classify each condition by normative force before relying on it.\end{CornellReferenceSummaryBox}
\Needspace{8\baselineskip}
\subsection*{8.5 Selection statements}
\begin{longtable}{@{}>{\RaggedRight\arraybackslash\columncolor{CornellReferenceCueGray}}p{0.28\textwidth}>{\RaggedRight\arraybackslash}p{0.64\textwidth}@{}}
\rowcolor{CornellReferenceNavy}\color{white}\textbf{Cue / recall prompt} & \color{white}\textbf{Notes}\\\endfirsthead
\rowcolor{CornellReferenceNavy}\color{white}\textbf{Cue / recall prompt} & \color{white}\textbf{Notes (continued)}\\\endhead
\arrayrulecolor{CornellReferenceLineGray}
\textbf{What rule applies to General?}\\[-1mm]{\scriptsize\CornellClauseRef{8.5.1}} & Frames the terminology, applicability, and relationships needed to interpret General; detailed rules remain in the following subclauses.\\\hline
\textbf{What rule applies to The if statement?}\\[-1mm]{\scriptsize\CornellClauseRef{8.5.2}} & Defines the core-language rules for The if statement, including formation, semantic effect, interactions with type and lookup, and any ill-formed or implementation-dependent cases.\\\hline
\textbf{What rule applies to The switch statement?}\\[-1mm]{\scriptsize\CornellClauseRef{8.5.3}} & Defines the core-language rules for The switch statement, including formation, semantic effect, interactions with type and lookup, and any ill-formed or implementation-dependent cases.\\\hline
\end{longtable}
\begin{CornellReferenceSummaryBox}Selection statements supplies the core-language semantics portion of this clause. The key review move is to connect its local wording to the clause-wide model, then classify each condition by normative force before relying on it.\end{CornellReferenceSummaryBox}
\Needspace{8\baselineskip}
\subsection*{8.6 Iteration statements}
\begin{longtable}{@{}>{\RaggedRight\arraybackslash\columncolor{CornellReferenceCueGray}}p{0.28\textwidth}>{\RaggedRight\arraybackslash}p{0.64\textwidth}@{}}
\rowcolor{CornellReferenceNavy}\color{white}\textbf{Cue / recall prompt} & \color{white}\textbf{Notes}\\\endfirsthead
\rowcolor{CornellReferenceNavy}\color{white}\textbf{Cue / recall prompt} & \color{white}\textbf{Notes (continued)}\\\endhead
\arrayrulecolor{CornellReferenceLineGray}
\textbf{What rule applies to General?}\\[-1mm]{\scriptsize\CornellClauseRef{8.6.1}} & Frames the terminology, applicability, and relationships needed to interpret General; detailed rules remain in the following subclauses.\\\hline
\textbf{What rule applies to The while statement?}\\[-1mm]{\scriptsize\CornellClauseRef{8.6.2}} & Defines the core-language rules for The while statement, including formation, semantic effect, interactions with type and lookup, and any ill-formed or implementation-dependent cases.\\\hline
\textbf{What rule applies to The do statement?}\\[-1mm]{\scriptsize\CornellClauseRef{8.6.3}} & Defines the core-language rules for The do statement, including formation, semantic effect, interactions with type and lookup, and any ill-formed or implementation-dependent cases.\\\hline
\textbf{What rule applies to The for statement?}\\[-1mm]{\scriptsize\CornellClauseRef{8.6.4}} & Defines the core-language rules for The for statement, including formation, semantic effect, interactions with type and lookup, and any ill-formed or implementation-dependent cases.\\\hline
\textbf{What rule applies to The range-based for statement?}\\[-1mm]{\scriptsize\CornellClauseRef{8.6.5}} & Explains the traversal or view contract for The range-based for statement; prove domain validity, lifetime, sentinel compatibility, and any borrowing or invalidation property.\\\hline
\end{longtable}
\begin{CornellReferenceSummaryBox}Iteration statements supplies the core-language semantics portion of this clause. The key review move is to connect its local wording to the clause-wide model, then classify each condition by normative force before relying on it.\end{CornellReferenceSummaryBox}
\Needspace{8\baselineskip}
\subsection*{8.7 Jump statements}
\begin{longtable}{@{}>{\RaggedRight\arraybackslash\columncolor{CornellReferenceCueGray}}p{0.28\textwidth}>{\RaggedRight\arraybackslash}p{0.64\textwidth}@{}}
\rowcolor{CornellReferenceNavy}\color{white}\textbf{Cue / recall prompt} & \color{white}\textbf{Notes}\\\endfirsthead
\rowcolor{CornellReferenceNavy}\color{white}\textbf{Cue / recall prompt} & \color{white}\textbf{Notes (continued)}\\\endhead
\arrayrulecolor{CornellReferenceLineGray}
\textbf{What rule applies to General?}\\[-1mm]{\scriptsize\CornellClauseRef{8.7.1}} & Frames the terminology, applicability, and relationships needed to interpret General; detailed rules remain in the following subclauses.\\\hline
\textbf{What rule applies to The break statement?}\\[-1mm]{\scriptsize\CornellClauseRef{8.7.2}} & Defines the core-language rules for The break statement, including formation, semantic effect, interactions with type and lookup, and any ill-formed or implementation-dependent cases.\\\hline
\textbf{What rule applies to The continue statement?}\\[-1mm]{\scriptsize\CornellClauseRef{8.7.3}} & Defines the core-language rules for The continue statement, including formation, semantic effect, interactions with type and lookup, and any ill-formed or implementation-dependent cases.\\\hline
\textbf{What rule applies to The return statement?}\\[-1mm]{\scriptsize\CornellClauseRef{8.7.4}} & Defines the core-language rules for The return statement, including formation, semantic effect, interactions with type and lookup, and any ill-formed or implementation-dependent cases.\\\hline
\textbf{What rule applies to The co\_\allowbreak{}return statement?}\\[-1mm]{\scriptsize\CornellClauseRef{8.7.5}} & Defines the core-language rules for The co\_\allowbreak{}return statement, including formation, semantic effect, interactions with type and lookup, and any ill-formed or implementation-dependent cases.\\\hline
\textbf{What rule applies to The goto statement?}\\[-1mm]{\scriptsize\CornellClauseRef{8.7.6}} & Defines the core-language rules for The goto statement, including formation, semantic effect, interactions with type and lookup, and any ill-formed or implementation-dependent cases.\\\hline
\end{longtable}
\begin{CornellReferenceSummaryBox}Jump statements supplies the core-language semantics portion of this clause. The key review move is to connect its local wording to the clause-wide model, then classify each condition by normative force before relying on it.\end{CornellReferenceSummaryBox}
\Needspace{8\baselineskip}
\subsection*{8.8 Declaration statement}
\begin{longtable}{@{}>{\RaggedRight\arraybackslash\columncolor{CornellReferenceCueGray}}p{0.28\textwidth}>{\RaggedRight\arraybackslash}p{0.64\textwidth}@{}}
\rowcolor{CornellReferenceNavy}\color{white}\textbf{Cue / recall prompt} & \color{white}\textbf{Notes}\\\endfirsthead
\rowcolor{CornellReferenceNavy}\color{white}\textbf{Cue / recall prompt} & \color{white}\textbf{Notes (continued)}\\\endhead
\arrayrulecolor{CornellReferenceLineGray}
\textbf{What rule applies to Declaration statement?}\\[-1mm]{\scriptsize\CornellClauseRef{8.8}} & Defines the core-language rules for Declaration statement, including formation, semantic effect, interactions with type and lookup, and any ill-formed or implementation-dependent cases.\\\hline
\end{longtable}
\begin{CornellReferenceSummaryBox}Declaration statement supplies the core-language semantics portion of this clause. The key review move is to connect its local wording to the clause-wide model, then classify each condition by normative force before relying on it.\end{CornellReferenceSummaryBox}
\Needspace{8\baselineskip}
\subsection*{8.9 Ambiguity resolution}
\begin{longtable}{@{}>{\RaggedRight\arraybackslash\columncolor{CornellReferenceCueGray}}p{0.28\textwidth}>{\RaggedRight\arraybackslash}p{0.64\textwidth}@{}}
\rowcolor{CornellReferenceNavy}\color{white}\textbf{Cue / recall prompt} & \color{white}\textbf{Notes}\\\endfirsthead
\rowcolor{CornellReferenceNavy}\color{white}\textbf{Cue / recall prompt} & \color{white}\textbf{Notes (continued)}\\\endhead
\arrayrulecolor{CornellReferenceLineGray}
\textbf{What rule applies to Ambiguity resolution?}\\[-1mm]{\scriptsize\CornellClauseRef{8.9}} & Defines the core-language rules for Ambiguity resolution, including formation, semantic effect, interactions with type and lookup, and any ill-formed or implementation-dependent cases.\\\hline
\end{longtable}
\begin{CornellReferenceSummaryBox}Ambiguity resolution supplies the core-language semantics portion of this clause. The key review move is to connect its local wording to the clause-wide model, then classify each condition by normative force before relying on it.\end{CornellReferenceSummaryBox}
\section{Language-Lawyer and Library-Usage Pitfalls}
\begin{CornellReferencePitfallBox}\begin{itemize}[label=\textcolor{CornellReferenceAmber}{$\blacktriangleright$}]
\item Ignoring scope and destruction when control leaves a block.
\item Assuming a range-based loop does not create hidden bindings.
\item Using a jump across declarations whose initialization rules prohibit the transfer.
\item Misreading a declaration/expression ambiguity.
\item Treating a note, example, or recommended practice as though it imposed a mandatory program rule.
\end{itemize}\end{CornellReferencePitfallBox}
\section{Programmer and Implementation Checklist}
\begin{CornellReferenceChecklistBox}\begin{itemize}[label=$\square$]
\item Build the control-flow graph including exceptional exits.
\item Track scope entry, initialization, and destruction on each path.
\item Check condition conversions and switch case constraints.
\item Expand range-based iteration conceptually before judging lifetime.
\item Validate return or coroutine-return rules against the function form.
\item For \CornellClauseRef{8.1}, verify preamble against the precise rule category and all cited dependencies.
\item For \CornellClauseRef{8.2}, verify label against the precise rule category and all cited dependencies.
\item For \CornellClauseRef{8.3}, verify expression statement against the precise rule category and all cited dependencies.
\item For \CornellClauseRef{8.4}, verify compound statement or block against the precise rule category and all cited dependencies.
\end{itemize}\end{CornellReferenceChecklistBox}
\section{Clause Synthesis}
\begin{CornellReferenceOrientationBox}Defines labels, blocks, selection, iteration, jumps, declaration statements, coroutine returns, and grammar ambiguity resolution. Organizes expressions and declarations into structured control flow. A sound review distinguishes syntax and participation from runtime preconditions, keeps notes and examples non-mandatory, and follows cross-clause dependencies before making a portability claim.\end{CornellReferenceOrientationBox}
\section{Self-Test Questions}
\begin{enumerate}
\item What role does Preamble play in the clause's overall model?
\item Which normative distinctions must be kept separate when applying Label?
\item How would you audit a program or implementation that depends on Expression statement?
\item Compare Compound statement or block with Selection statements; what different question does each answer?
\item What role does Selection statements play in the clause's overall model?
\item Which normative distinctions must be kept separate when applying Iteration statements?
\item How would you audit a program or implementation that depends on Jump statements?
\item Compare Declaration statement with Ambiguity resolution; what different question does each answer?
\item Scenario: a program relies on an unstated implementation behavior while using Ambiguity resolution. What must be checked before calling the program portable?
\item Scenario: a program relies on an unstated implementation behavior while using General. What must be checked before calling the program portable?
\end{enumerate}
\clearpage\section{Answer Key}
\begin{enumerate}
\item Preamble is one part of the clause's core-language semantics model. Its role is understood by connecting the local rule in 8.1 to the clause orientation and its dependent concepts.
\item Keep formation and well-formedness separate from runtime preconditions and effects. Also distinguish mandatory wording from notes, implementation choice from undefined behavior, and interface declarations from detailed contracts; see 8.2.
\item Audit Expression statement by locating the controlling wording in 8.3, listing inputs and type constraints, proving any preconditions, recording effects and failure behavior, and checking lifetime, invalidation, complexity, and synchronization where applicable.
\item Compound statement or block and Selection statements occupy different steps or facility groups. Analyze each under its own subclause, then follow the explicit dependency between them rather than transferring one set of guarantees to the other; begin at 8.4.
\item Selection statements is one part of the clause's core-language semantics model. Its role is understood by connecting the local rule in 8.5 to the clause orientation and its dependent concepts.
\item Keep formation and well-formedness separate from runtime preconditions and effects. Also distinguish mandatory wording from notes, implementation choice from undefined behavior, and interface declarations from detailed contracts; see 8.6.
\item Audit Jump statements by locating the controlling wording in 8.7, listing inputs and type constraints, proving any preconditions, recording effects and failure behavior, and checking lifetime, invalidation, complexity, and synchronization where applicable.
\item Declaration statement and Ambiguity resolution occupy different steps or facility groups. Analyze each under its own subclause, then follow the explicit dependency between them rather than transferring one set of guarantees to the other; begin at 8.8.
\item First identify the exact requirement in 8.9, then classify the relied-on behavior as required, unspecified, implementation-defined, conditionally supported, or undefined. Portability requires satisfying the program obligations and avoiding dependence on a choice not guaranteed in the target environments.
\item First identify the exact requirement in 8.7.1, then classify the relied-on behavior as required, unspecified, implementation-defined, conditionally supported, or undefined. Portability requires satisfying the program obligations and avoiding dependence on a choice not guaranteed in the target environments.
\end{enumerate}
\section{Key-Term Recap}
\begin{longtable}{@{}>{\RaggedRight\arraybackslash}p{0.28\textwidth}>{\RaggedRight\arraybackslash}p{0.64\textwidth}@{}}\toprule\textbf{Term} & \textbf{Working meaning}\\\midrule\endfirsthead\toprule\textbf{Term} & \textbf{Working meaning (continued)}\\\midrule\endhead\bottomrule\endfoot
statement & A syntactic unit controlling execution or introducing a declaration.\\
block & A compound statement that groups statements and forms a scope.\\
selection statement & A conditional control structure choosing among branches.\\
iteration statement & A control structure that repeats a substatement.\\
jump statement & A statement that transfers control, returns, or produces a coroutine result.\\
condition & The declaration or expression whose converted value controls selection or iteration.\\
\end{longtable}
\CornellTakeaway{Statements are not only control flow: they also determine scopes, initialization points, destruction, and permitted transfers.}
\end{document}
